\clearpage
\setcounter{page}{1}
\renewcommand{\thepage}{S\arabic{page}}
\maketitlesupplementary


\begin{appendix}

\section*{Overview}

This supplementary material provides the construction and screening details of
DeformNeXt-10k, the implementation of semantic correspondence and
matching-aware feature aggregation, additional training and evaluation
protocols, and qualitative results.  We report all evaluation variants with
their input and alignment assumptions explicitly stated; ground-truth
correspondence is used only to construct training supervision and is never
provided to the model at inference time.


\section{DeformNeXt-10k Construction Details}
\label{sec:supp_deformnext}

\paragraph{Source assets and canonical correspondence.}
We construct DeformNeXt-10k from textured PartNeXt assets and their hierarchical part annotations~\cite{wang2026partnext}.
For each source mesh $M_0=(V_0,F,A_0)$, we generate edited variants $M_{1,k}=(V_{1,k},F,A_0)$, modifying vertex positions while preserving the ordered face array and source UV, material, and corner attributes.
Consequently, corresponding points can be identified by the same face index and barycentric coordinates, without nearest-neighbor reassignment.
This correspondence is used to construct training labels for cross-state matching; it is not assumed for surfaces decoded by the generative model.

\paragraph{Donor-guided deformation.}
We select same-category donor objects using fixed semantic-role and structural-compatibility heuristics.
Roles are grouped semantically rather than requiring one-to-one correspondences between individual mesh components, allowing objects with different part counts and disconnected components.
Eight renders of each donor with its original appearance are used to personalize a pretrained SDXL prior through GeoStyle~\cite{choi2026image}.
Donor-specific DreamBooth LoRA is trained for $800$ updates (rank $16$, resolution $512$, learning rate $10^{-4}$), and its weights can be reused across source--donor pairs.
We use the personalization caption ``A photo of TOK sculpture'' and the deformation prompt ``A TOK style sculpture of a \texttt{<category>}''.
The donor supplies a morphology prior for proposing edits; neither the donor geometry nor its images are provided to the appearance-retargeting model at inference.

\paragraph{Part-aware assembly.}
GeoStyle proposes per-face Jacobian changes through image-space score distillation.
To better preserve the source assembly, we deform selected primary surfaces while treating complementary physical components as rigid followers.
Hierarchical semantic labels identify candidate roles, but do not provide ground-truth mechanical joints.
We instead derive physical components from source connectivity and coincident-vertex classes, and estimate attachment constraints from nearby source surfaces using vertex-to-triangle patches, barycentric coordinates, and rest-frame offsets.
The default attachment radius is $0.006$ times the source diagonal, with at most $96$ contact samples per component.
Reconstruction alternates four projected Poisson/rigid steps, followed by a Dirichlet solve for follower positions.
These constraints improve assembly coherence without guaranteeing collision-free geometry.

\paragraph{Optimization and optional refinement.}
The main production configuration performs $300$ deformation iterations at resolution $512$ with batch size one (seed $37$), using main/auxiliary score-distillation weights of $2.5/0.25$ and a cage-loss epoch of $83$.
This shorter budget was chosen for dataset-scale throughput; it is not assumed equivalent to the reference $3{,}600$-iteration GeoStyle configuration (resolution $1024$, batch size four).
The final snapshot also retains a small number of earlier longer-budget variants with their original provenance.
We additionally consider an optional balanced Jacobian--normal--Laplacian (JN-L) refinement of a coarse deformation proposal, optimized for $120$ Adam iterations at learning rate $0.003$.
The refinement combines Jacobian consistency, source-crease-aware normal preservation, and an area-weighted residual umbrella Laplacian.
For coarse vertex positions $\bar V$ and residual displacements $r_i=v_i-\bar v_i$, the latter term is
\begin{equation}
    \mathcal{L}_{\mathrm{L}} =
    \sum_i m_i\left\|r_i-\sum_{j\in\mathcal{N}(i)}w_{ij}r_j\right\|_2^2,
    \label{eq:supp_deformnext_laplacian}
\end{equation}
where $m_i$ is an area weight and $w_{ij}$ denotes row-normalized, inverse-edge-length weights on the source mesh.
Boundary, nonmanifold, and sharp-crease edges are excluded from this neighborhood.
The Laplacian weight is fixed from the initial gradient-norm ratio to the Jacobian and normal terms, rather than tuned using final evaluation scores.
Both the original part-aware (PA) proposal and its refined version are considered independently.

\paragraph{Screening and paired observations.}
Each candidate is assessed against the original source $M_0$, rather than against its donor or initial proposal.
We check canonical correspondence and finite geometry, meaningful deformation, severe assembly failures, paired RGB usability, survival under native 3D representations, and compatibility with the training loader.
Multi-view, model-assisted screening also identifies conspicuous spikes, detached components, and assembly collapse; it is not equivalent to exhaustive human review or a guarantee of intersection-free meshes.
All accepted pairs retain the source-defined coordinate frame (extent $0.48$), original source materials, and matched calibrated cameras, without independently normalizing the edited geometry.
For retargeting, an edited $M_1$ target-view render specifies the desired state, whereas undeformed $M_0$ source-view renders supply appearance evidence; neither donor imagery nor ground-truth correspondence is passed to the generator at inference.
Approved distinct variants are retained, while near-duplicate final geometries are merged.
The frozen snapshot contains $5{,}611$ PA, $4{,}548$ balanced JN-L, and six earlier bootstrap variants.

\paragraph{Dataset scope and splits.}
The resulting dataset contains $10{,}165$ accepted edits from $3{,}483$ source assets spanning $50$ categories.
A fixed approximately $90/5/5$ train/validation/test split contains $9{,}149/508/508$ edits, with all variants of each source UID and existing source-geometry cluster assigned to a single split.
The split is disjoint in source identities and geometry clusters, but donor identities are not constrained to be split-disjoint; overlap with earlier multi-view pretraining data has not been established.
The dataset is intended to provide plausible and observable shape edits with source-material correspondence, rather than physically certified assemblies or an estimate of the success rate of the underlying deformation procedure.




































\section{Semantic Feature Matching Details}
\label{sec:supp_semantic_matching}

\paragraph{Descriptors and candidate retrieval.}
For each observed target patch $m$, we retrieve the top-$64$ candidate patches $\mathcal{C}_m$ across the three undeformed source views using cosine similarity of semantic descriptors.
We use $1024$-dimensional layer-12 FRAME descriptors from a frozen VGGT-$\Omega$ backbone~\cite{wang2026vggtomega} for correspondence estimation.
These matching descriptors are distinct from the native DINOv3 features subsequently aggregated for 3D generation (Sec.~\ref{sec:supp_matching_aggregation}).
Candidate retrieval defines the support over which matching scores are refined; the decoder cannot recover a ground-truth match absent from this initial proposal set.

\paragraph{Matching decoder.}
We use a lightweight D1 decoder to refine the candidate ranking under geometric deformation.
It contains one query-only cross-attention block with four heads of width $32$ and a $128$-dimensional feed-forward bottleneck.
The decoder contextualizes the target query, while the retrieved candidate descriptors remain unchanged.
A LayerNorm--MLP scalar head with hidden dimension $64$ adds a bounded $0.25\tanh$ residual to the contextualized cosine score.
Thus, the final score $s_{mk}$ combines descriptor similarity with a learned, candidate-specific correction; the refined scores are used to construct the matching probabilities in Sec.~3.4.

\paragraph{Correspondence supervision.}
As described in Sec.~\ref{sec:supp_deformnext}, each source--target mesh pair preserves its ordered face indices and barycentric surface correspondence.
We use these correspondences together with training-only depth checks to identify valid positive source patches $\mathcal{P}_m$ for each target patch.
The multi-positive listwise objective in Sec.~3.3 assigns probability to the intersection $\mathcal{C}_m\cap\mathcal{P}_m$, allowing more than one valid positive within the retrieved support.
Only queries with at least one valid positive candidate contribute to the loss; losses are averaged over eligible queries within each asset, and assets without eligible queries do not update the matcher.
Queries without an observable correspondence and cases where the proposal set misses all positives are recorded separately rather than treated as reliable negatives.
Ground-truth surface correspondence and depth checks are used for supervision and diagnostic scoring, not for inference-time source selection or visibility gating.

\paragraph{Training and freezing.}
Of the $9{,}149$ geometry-stage training examples, $9{,}143$ are certified for matcher training; six teacher-rejected cases are excluded from this step only.
We train the D1 decoder for seven passes ($62{,}545$ updates) using AdamW with learning rate $10^{-4}$, weight decay $0.01$, and gradient clipping at $1$.
The VGGT-$\Omega$ backbone is frozen during decoder training, and the complete matcher is frozen during subsequent 3D generative fine-tuning.
No matching-head training or calibration is performed on the test split.

\paragraph{Matching diagnostics.}
On $50{,}493$ observable source-camera-ray/target-patch relations from $32$ development source clusters, raw cosine matching on the mean-ray support achieves $1.25\%$ exact Top-1 accuracy.
Contextualized decoder similarity increases this to $7.73\%$, and scalar score refinement further increases it to $8.56\%$.
Exact Top-8 coverage increases from $11.09\%$ to $28.47\%$ and $29.92\%$, respectively.
Exact Top-64 coverage remains $50.76\%$ for these variants because they rerank the same proposal support; the additional coverage from combining observed target-member proposals is analyzed separately in Sec.~\ref{sec:supp_matching_aggregation}.
These retrieval diagnostics measure correspondence ranking and candidate coverage, not downstream 3D generation quality.










































\section{Matching-Aware Feature Aggregation Details}
\label{sec:supp_matching_aggregation}

\paragraph{Ray grouping and observation support.}
The geometry (SS) stage uses a $16^3$ conditioning grid and $32\times32$ image-patch maps; SS64 denotes the decoded occupancy resolution rather than the conditioning-grid size.
For each source camera $j$ and source patch $p$, we group conditioning sites that project to the same source-camera patch address into a ray $r=(j,p)$.
The corresponding set $\mathcal{M}_r$ contains the \emph{distinct observed target patches} associated with those grid sites.
Repeated projections to the same target patch contribute only once, avoiding an implicit preference for patches intersected by more grid sites.
Ray membership and observation support follow the calibrated cameras and image alpha; neither target surface depth nor ground-truth visibility is required at inference.

\paragraph{Ray-level matching posterior.}
As described in Sec.~3.3, each observed target patch $m$ retrieves up to $64$ candidates $\mathcal{C}_m$ from the three source images, using frozen semantic matching descriptors and the D1 decoder's refined scores $s_{mk}$.
We normalize these scores over $\mathcal{C}_m$ with temperature $\tau=0.05$.
A source candidate $k$ indexes both its source view and image patch; candidates absent from a member's proposal set receive zero probability for that member.
We average the normalized distributions across $\mathcal{M}_r$, retain the $64$ highest-probability candidates from the resulting ray-level mixture, and renormalize their weights, as defined in Eqs.~\eqref{eq:matching_probability}--\eqref{eq:ray_matching}.
Unlike matching a single averaged target descriptor, this procedure retains alternative correspondence hypotheses contributed by different target patches before truncation.

\paragraph{Aggregation of native source features.}
The matching descriptors and the values transferred to the generator are distinct.
Matching uses $1024$-dimensional layer-12 FRAME descriptors from the frozen VGGT-$\Omega$ backbone, whereas the geometry stage gathers $1024$-dimensional \emph{native DINOv3 features} at the selected source patches.
The final aggregation is the weighted sum in Eq.~\eqref{eq:matching_aggregation}, using the ray-level matching probabilities as weights.
The selected Fixed Top-K design uses these native values directly, without an additional learned query/key attention module, value projection, or output projection.
Consequently, all conditioning sites in a source-camera ray receive the same retrieved evidence, which is subsequently used to correct the source contribution of the native projection conditioner (Sec.~3.5).

\paragraph{Unsupported observations and fallback.}
Rays with no observed target members or insufficient valid proposals abstain from cross-state retrieval, leaving the original source-conditioning path unchanged.
Importantly, image alpha and camera projection specify observation support rather than physical source--target co-visibility.
Likewise, a concentrated matching posterior is not necessarily a calibrated estimate of correspondence correctness.
The fallback therefore handles absent evidence but cannot eliminate confidently incorrect matches.
Ground-truth surface correspondences and depth checks are used for matcher supervision or diagnostics only, never as inference-time visibility masks.

\paragraph{Aggregation alternatives and diagnostics.}
Our comparisons include retrieval from the actual Top-1 candidate and uniform weighting over the \emph{same} retained Top-K support.
The latter isolates the role of matching-score weights, whereas Top-1 changes both the number of retrieved candidates and their aggregation.
We also evaluate learned score refinement as an alternative to fixed weighting.
To preserve the initial matching distribution, its per-head attention weights are anchored to the frozen prior:
\begin{equation}
    a_{rhk}(\theta)=
    \operatorname{softmax}_{k\in\mathcal{K}_r}
    \left[\log p_r(k)+s_{rhk}(\theta)-s_{rhk}(\theta_0)\right],
    \label{eq:supp_anchored_matching_attention}
\end{equation}
where $h$ indexes four heads and $\theta_0$ denotes initialization.
The head probabilities are averaged before aggregating native values; at $\theta=\theta_0$, the weights recover the Fixed posterior.
This learned refinement is an evaluated alternative, not part of the selected Fixed Top-K pathway.

On $50{,}493$ observable relations from $32$ development source clusters, combining target-member proposals increases exact Top-64 matching coverage from $50.76\%$ to $73.36\%$.
Among $31{,}408$ observable relations where Top-1 misses but Top-K contains a ground-truth candidate, the mean probability assigned to valid candidates is $1.63\%$ for post-hoc uniform weighting, $6.46\%$ for Fixed weighting, and $6.18\%$ for learned refinement.
These quantities characterize candidate availability and probability allocation, rather than downstream generation quality; the post-hoc uniform diagnostic is distinct from the separately trained uniform-weighting control.

























\section{Target-Aware Feature Retargeting Details}
\label{sec:supp_target_retargeting}

\paragraph{Retargeting within the native conditioner.}
The geometry stage retains the pretrained multi-view projection conditioner, including its target-view contribution and global conditioning tokens.
For a source-camera ray $r$, let $\mathbf{E}_r$ be the matching-weighted source evidence from Sec.~3.4 and $\mathbf{S}_r$ the native source feature already represented in the conditioning pathway.
Rather than injecting $\mathbf{E}_r$ as an independent feature, we use the difference $\mathbf{E}_r-\mathbf{S}_r$ to express a correction relative to the existing source contribution.
This avoids treating the retrieved evidence as an additional observation and retains the scale and structure of the pretrained conditioning pathway.
The correspondence decision is shared by conditioning-grid sites belonging to the same source-camera ray, as described in Sec.~\ref{sec:supp_matching_aggregation}.

\paragraph{Target context and adaptive source use.}
The retargeter combines the source-feature difference with a native target context $\mathbf{T}_r$ pooled over the conditioning-grid sites of ray $r$.
The target context is distinct from the set of observed target patches used to construct the matching posterior: the latter determines which source candidates are retrieved, whereas $\mathbf{T}_r$ provides target-state information for adapting their features.
A learned source-use probability $u_r$ controls the contribution of the correction, together with a bounded signed gate $g$ on the direct feature difference.
The learnable retargeter $\mathcal{R}_{\phi}$ produces an additional feature residual conditioned on $[\mathbf{E}_r-\mathbf{S}_r,\mathbf{T}_r]$, yielding the correction in Eq.~\eqref{eq:feature_retargeting}.
The source-use probability is a conditioning weight, not a calibrated estimate of physical surface visibility or matching correctness.

\paragraph{Residual injection and initialization.}
Source corrections are organized in their respective source-view patch maps and bilinearly projected into the 3D conditioning grid using the original calibrated camera operators.
Their sum is normalized by the total number of conditioning views, consistent with the native projection-based aggregation.
The target-view path and global conditioning tokens are not replaced by the retargeted source evidence.
We initialize the signed gate to zero and the retargeter output to zero, so the correction initially vanishes and the pretrained conditioning function is recovered.
If a ray has no usable matching support, its correction is omitted and the native source contribution is retained rather than replaced with a zero feature.
This fallback prevents unsupported updates, although it does not detect confidently incorrect correspondences.
Preserving the target conditioning path does not guarantee pixel-exact reproduction of the edited target image after 3D generation.

\paragraph{Stage-specific source adaptation.}
We adapt the geometry and appearance stages separately rather than imposing one correspondence aggregation strategy on the complete cascade.
The sparse-structure stage (SS64) uses the Fixed Top-K evidence of Sec.~3.4 followed by the gated residual retargeter.
The intermediate Shape512 and Shape1024 stages remain frozen.
Appearance synthesis (PBR1024) operates on the generated geometry and uses its own source-conditioning pathway with a separately trained Top-1 retargeter and native DINOv3/NAF features.
The main comparison retains this pre-existing Top-1 PBR checkpoint and does not replace it after inspecting the test set.
In the common-501 fixed-geometry control, original conditioning, Top-1, and Top-K have comparable source-fidelity metrics: Top-K improves LPIPS relative to Top-1, but its DreamSim and CLIP intervals include zero, and none of its three intervals against original conditioning excludes zero.
Accordingly, we do not claim that Top-1 is uniformly preferable for appearance or that appearance retargeting provides a broad gain.

\paragraph{Fine-tuning and controlled comparisons.}
We continue the supplied multi-view base-plus-LoRA checkpoints, using the original SS64/Shape512 checkpoint at step $4{,}210$ and Shape1024/PBR1024 checkpoint at step $16{,}850$.
Image encoders, the VGGT-$\Omega$ matching backbone and decoder, and the intermediate shape stages remain frozen during generative fine-tuning.
The geometry stage is trained for $480$ updates with global batch size $96$, LoRA rank/alpha $16/16$, and denoiser/conditioner learning rates $10^{-5}/10^{-4}$.
The appearance stage uses $1{,}900$ updates, global batch size $24$, and denoiser/conditioner learning rates $5\times10^{-6}/10^{-4}$.
Geometry-stage alternatives are compared with common frozen downstream shape models and the same Top-1-retargeting PBR model; appearance-stage alternatives are compared on identical generated geometry, guides, texture noise, rendering triangles, and references.
Modules removed in ablations are physically omitted rather than replaced with zero-filled inputs.
These controls separate changes in geometry conditioning from changes in appearance synthesis.

\paragraph{Common-success evaluation protocol.}
The cross-method comparison uses the 501 test edits successfully produced by
every evaluated method, spanning 133 source-geometry clusters.  Generation
never receives a scale measured from the edited target mesh.  Methods that
natively accept calibrated cameras retain that interface; the original
single-view Pixal3D baseline additionally appears at its native 1536
resolution with estimated camera parameters.  For every method, we report
the native output and a source-only Sim(3) variant whose transform is fitted
to the input source mesh $M_0$.  The edited target mesh $M_1$ is opened only
after this transform has been fixed and is used solely for CD and F-Score
evaluation.  Source preservation uses cameras 9 and 10, target-image metrics
use camera 8, and all images are rendered by the same Cycles protocol.

The \emph{MV-Pixal3D (zero-shot)} control uses the pretrained mixed385
multi-view checkpoints without DeformNeXt optimization and receives the same
edited target plus three original source RGB images as our method.  The
\emph{Pixal3D (target 1-view)} control instead receives only the edited target
RGB and its calibrated camera in the 1024 cascade.  Keeping these controls
separate from the native-1536 Pixal3D row exposes the effects of multi-view
input, camera access, and task-specific fine-tuning without treating any of
them as an isolated proposed-module contribution.













\end{appendix}
